\section{Results}

We now summarize the empirical behavior of the 16-D implementation on the
RTX~A6000.  For each $n_{\text{train}}$ in $\{4\text{k}, 8\text{k}, 16\text{k},
32\text{k}\}$ we set $n_{\text{test}} = n_{\text{train}}/8$, draw data from a
simple 16-D Gaussian mixture, and run three baselines: scikit-learn KDE,
Torch SD-KDE (GEMM-based), and Triton SD-KDE (Tensor-Core GEMMs for both KDE
and score).  Figure~\ref{fig:nd_runtime} shows that the Triton SD-KDE rapidly
outpaces both baselines as $n$ grows, while remaining numerically close to the
Torch SD-KDE reference.

We also measure utilization by combining the flop model above with the
measured runtimes.  As shown in Figure~\ref{fig:triton_sd_kde_nd_util}, the
Torch SD-KDE baseline reaches only a few percent of the A6000 Tensor Core
peak, whereas the Triton implementation climbs into the multi-digit range once
$n_{\text{train}}$ exceeds $8\text{k}$.  This confirms that the 16-D
Tensor-Core formulation is firmly compute-bound and that additional tuning
effort should focus on kernel fusion and occupancy rather than memory traffic.

To compare against a specialized KDE implementation, we also benchmarked PyKeOps
LazyTensor operators at $n_{\text{train}}=32\text{k}$, $n_{\text{test}}=4\text{k}$.
PyKeOps currently represents the state of the art for large-scale kernel
reductions, so we report both its Gaussian KDE runtime and its SD-KDE runtime as
a reference point.  Table~\ref{tab:pykeops} shows that Triton SD-KDE runs
$1.4\times$ faster than the PyKeOps 16-D KDE baseline and $9.8\times$ faster
than the PyKeOps 16-D SD-KDE implementation.

\begin{table}[h]
  \centering
  \small
  \setlength{\tabcolsep}{4pt}
  \begin{tabular}{lcc}
    \toprule
    Method & Runtime (ms) & Rel to Triton \\
    \midrule
    Triton 16-D SD-KDE & $2.23$ & $1\times$ \\
    PyKeOps 16-D KDE   & $3.06$ & $1.4\times$ \\
    PyKeOps 16-D SD-KDE& $21.73$ & $9.77\times$ \\
    \bottomrule
  \end{tabular}
  \caption{Runtime comparison at $n_{\text{train}}=32\text{k}$,
           $n_{\text{test}}=4\text{k}$.  PyKeOps rows report Gaussian KDE and
           SD-KDE (LazyTensor) runtimes, while Triton executes the full SD-KDE
           pipeline.}
  \label{tab:pykeops}
\end{table}

\begin{figure*}[t]
  \centering
  \includegraphics[width=\textwidth]{figures/nd-runtime.pdf}
  \caption{Runtime comparison for 16-D KDE/SD-KDE across $n_{\text{train}}$
           up to $32{,}768$ ($n_{\text{test}} = n_{\text{train}}/8$).}
  \label{fig:nd_runtime}
\end{figure*}
